\section{Full cut-boundary cells and source-certified piece profiles}
\label{sec:cut-boundary}

The hierarchy must retain actual geometric attachments after a normal cut.
The preceding chamber and prism operators locate face-connected pieces and
their interface sets, but they do not reconstruct the boundary of the cut
manifold. We now construct its complete two-dimensional cell structure in
binary families, with cyclic polygon attachments, edge endpoints, both sides
of normal intersection vertices and a carrier into the original cut chambers.
The resulting operator determines all boundary genera and each cut piece's
boundary count and total boundary genus without expanding normal sheets.

This is the boundary of $M\setminus\operatorname{int}\nu(F)$ for the supplied
normal surface $F$, including its one-sided components. It is not the base
of the full interior parallelity bundle. The ambient source contract remains
finite, compact, connected, orientable and with one torus boundary. No source
coordinate is divided by a common factor: repeated surfaces change the cut
manifold, not just the spelling of its input.

\subsection{Two horizontal sides and the original cut boundary faces}

Give each normal triangle and quadrilateral two local polygon copies. For a
triangle near vertex $v$, side zero faces $v$ and side one faces away from it.
For quad type $j$, side zero faces $A=\{0,j+1\}$ and side one faces its
complement. These are local sides, not an assumed global coorientation.
Number the two sides in contiguous blocks for each of the seven disc types.

On a source face $f$, the arc stack near corner $v$ has length
$a_v=t_v+q_{f,v}$. Its first $t_v$ arcs are triangles, followed by the
applicable quad arcs. Triangle rank increases from $v$. Quad rank increases
from $A$ and decreases from the other side. Consequently a face arc's rank
maps to its horizontal polygon by at most two affine pieces of slope $+1$
or $-1$. Each arc also has an inward and an outward side relative to its
face corner; this specifies which local horizontal copy it meets.

For a glued face, match equal arc ranks and equal corner sides through its
vertex permutation. Intersect the two rank partitions and retain both local
side maps. This includes self-gluings and the monodromy of one-sided surfaces.
For a boundary face, retain its central region and three corner arms. Arc
rank $r$ has inward region $r$, and outward region $r+1$, except that the
last outward region is the centre. Glue these regions to their corresponding
horizontal polygons.

Finally glue the cut segments of each original boundary edge between its
two incident boundary-face occurrences. If the edge has $w$ normal points,
there are $w+1$ gap segments. The face sends its first corner stack to one
arm, its middle gap to its centre, and its last stack to the other arm.
Global edge orientation comes from the validated signed edge union, including
loop edges. This partitions the gap ranks into three affine pieces on each
face. These joins are essential: horizontal-side matching alone would leave
the original boundary faces disconnected.

\begin{proposition}[Exact compressed cut-boundary complex]
The polygons and joins above are exactly the full boundary of the cut-open
manifold. Each edge has two polygon incidences and each vertex link is a
circle. Polygon attachments have constant-length cyclic recipes with affine
rank parameters. The construction has $O(T)$ cell families and interval
joins, with endpoint bit length $O(B+\log(T+1))$ for maximum coordinate bit
length $B$. No sheet, edge point or polygon family is expanded.
\end{proposition}
\begin{proof}
Inside each source tetrahedron the frontier of a normal regular neighbourhood
has the two stated local copies of every normal disc. Across a glued face,
each copy of an arc is attached to the copy on the same corner side; the
matching equations identify exactly its rank. This reconstructs the entire
horizontal frontier, with both local sides glued even when no consistent
global side labelling exists.

The remaining boundary comes from the source boundary faces cut along their
normal arcs. Their centre and corner-arm regions are discs. An inward arc
copy bounds its rank-$r$ region; an outward copy bounds the next region or
the centre. The cut original edge gaps glue the two adjacent source faces.
These are exactly the pieces and identifications in the boundary of
$M\setminus\operatorname{int}\nu(F)$.

At a normal intersection with a source edge there are two distinct cut
vertices, indexed by the side of the globally oriented edge. Rank from a
face corner maps to its edge point by a translation or reflection, and its
corner side determines the cut-vertex side. Original boundary vertices
remain unchanged. The validated manifold edge links and vertex links imply
that the complete boundary neighbourhoods are discs. In particular, these
are genuine boundary vertices, not a count inferred from prism 0/1-handles.

Horizontal triangle and quad cycles have three and four arcs. A corner-arm
region is triangular at rank zero and quadrilateral at positive rank. A
central region has between three and six edges. Cyclic recipes retain their
source face, corner, local side, oriented rank and the optional inner arc
at positive rank. All edge endpoints are source-labelled affine maps;
original gap endpoints split only at the first and last gap. Their cyclic
ordering therefore describes full polygon attachments, including repeated
vertices and self-identifications.

An internal face contributes at most 18 horizontal joins, a boundary face
at most 15 horizontal/original joins, and a boundary edge at most seven
gap joins. There are at most $2T$ glued face pairs, $4T$ boundary faces
and $6T$ boundary edges, giving at most $138T$ edge-family joins. There
are at most $30T$ polygon families and $16T$ vertex families. Every cell
count is bounded by a constant multiple of $T+D$, with
$D\le7T2^B$. All processing uses these binary counts and constantly many
events per source face or edge.
\end{proof}

The source orientation fixes opposite signed incidences on a shared edge.
The producer uses the standard triangle and quad cycles; independent replay
instead builds the crossing-edge polygon and orients it with integer midpoint
cross products in the unit tetrahedron. This avoids a shared orientation-table
assumption. Small literal cell audits check cyclic closure, signed incidences
and connected vertex links independently of orbit discovery.

\subsection{Owned vertices, edges and exact boundary genera}

Each polygon receives one face weight. Assign each normal-arc edge copy to
its canonical source horizontal polygon, and each original gap segment to
its canonical boundary-face occurrence, with weight minus one. Assign the
two cut vertices at a normal edge point through one canonical incident
tetrahedron's crossing discs. Assign each original boundary vertex through
one incident innermost arm or centre. Each cell is owned exactly once.

Use three coordinates: Euler weight, horizontal-polygon count and
original-polygon count. The latter two preserve source provenance and
distinguish otherwise equal genus profiles. Weighted orbit sums over polygon
adjacency give each connected boundary component's $V-E+F$, and group equal
profiles with a binary multiplicity. The closed boundary is orientable, so
\[
                   g=1-\chi/2.
\]
An odd Euler value or a value greater than two is rejected.

For a global conservation check let $D$ be the normal disc count, $A$ the
normal arc count over canonical source faces, $W$ the total normal edge
weight, and use subscripts $\partial$ for the original boundary. The cells
have counts
\[
\begin{split}
V'&=V_\partial+2W,\\
E'&=E_\partial+W_\partial+2A,\\
F'&=F_\partial+A_\partial+2D.
\end{split}
\]
Every boundary normal arc has two ends, and every boundary edge point is
incident to two such arcs; hence $A_\partial=W_\partial$.
Because the original boundary is a torus, subtraction gives
\[
                   V'-E'+F'=2(W-A+D)=2\chi(F).
\]
This checks global mass but does not replace the per-component orbit proof.
The independent checker computes horizontal Euler coefficients directly from
canonical edge owners and face owners, instead of repeating the producer's
individual edge contributions.

\subsection{Boundary ownership inside cut pieces}

Every horizontal polygon has a specific cut chamber on its local side.
A triangle's inward copy meets its rank-$r$ triangle-arm chamber. Its
outward copy meets the next arm chamber, except at the last triangle,
where it meets the appropriate quad-chain endpoint or the centre. A quad's
two sides meet chain chambers $r$ and $r+1$. The original boundary regions
use the face/chamber maps of Section~\ref{sec:cut-complement}.

These carriers are a partition of all boundary polygons into affine maps
to cut chambers. There are at most $42T$ carrier intervals. Join the
boundary polygon system to the existing chamber system through these maps.
The joint orbits are exactly the connected pieces of the cut manifold:
boundary joins lie within a genuine cut piece, and every polygon is attached
to its actual chamber. No ownership arrow crosses the cutting surface.

To count boundary components inside a cut piece, first obtain a compressed
least transversal of the boundary polygon orbits. Its original-source
intervals contain exactly one representative per boundary component.
Give these representatives one additional marker coordinate. Joint weighted
replay then obtains, for every cut piece, the boundary count $b$, boundary
Euler sum $\chi_\partial$, its horizontal and original polygon counts and
its number of source chambers. In particular
\[
              \sum_{i=1}^{b}g_i=b-\chi_\partial/2.
\]
This is a total genus, not an asserted distribution of individual genera
inside a piece with several boundary components. The complete boundary
cells and their carriers retain the underlying attaching data.

\begin{proposition}[Polynomial binary boundary and ownership query]
Full cut-boundary construction, its genus inventory and cut-piece boundary
counts and total genera have polynomial bit complexity in $T$ and $B$.
Independent verification additionally charges the supplied cell records
and all supplied trace events. Neither operation enumerates a represented
boundary component or a parallel surface copy.
\end{proposition}
\begin{proof}
The boundary interval universe has
$N_\partial=2D+F_\partial+A_\partial\le6D+4T$ points and at most $138T$
joins. The chamber universe has $T+D$ points and at most $20T$ joins.
Adding at most $42T$ carriers gives at most $200T$ joint joins and a universe
of at most $7D+5T$ points. Their endpoint bit length is
$O(B+\log(T+1))$.

Boundary discovery uses the existing monotone Agol--Hass--Thurston trace
and compressed least-transversal compiler. Weighted boundary replay reuses
that authenticated trace. The joint system requires one further discovery.
The fixed weight dimensions are three and five. Initial coefficients are
bounded owned-cell counts; intermediate and final totals have polynomial
binary lengths. The polynomial interval-orbit, weighted-transport and
transversal bounds of the maintained operators apply, with their complete
arithmetic and trace accounting\cite{aht-orbits}. Canonical family sorting
and independent cell reconstruction retain polynomial source/bit cost.
No multiplicity is used as a loop bound.
\end{proof}

\subsection{Source replay and the topology that remains}

The API is \code{normal\_cut\_boundary}. Independent checking uses
\code{verify\_normal\_cut\_boundary}. The schema is
\code{normal-cut-boundary-v1}. It binds the actual source and unscaled
coordinates, every face/edge/vertex family and oriented cycle, every chamber
carrier, the boundary transversal and both weighted proofs. Rank-event
reconstruction and midpoint orientation import no boundary producer.
Signed hexadecimal transport supports large binary fields.

Discovery caps share one allowance across boundary discovery and the joint
query; weighted trace reuse does not start another discovery. Incomplete
queries return no geometry, summary or certificate. The verifier's optional
operation cap charges all three supplied traces, including the repeated
boundary trace, and means unverified on exhaustion. Cooperative cancellation
propagates through source construction, compilation and replay.

Boundary spheres do not certify a ball. The retained solid-torus/RP$^3$
connected-sum source, cut along its standard vertex disc number six, has
one sphere boundary but first homology $\mathbb Z/2$. Regina's independent
cut confirms it is not a ball. The new API therefore returns no ball or
unknot verdict. Pattern essentiality and irreducibility would need their
own applicable certificates.

The constructor supplies actual boundary cells and their source carriers,
not the full compressed three-dimensional interior or an admissible handle
structure. It does not compute the interior parallelity base, collapse its
handles, transport an arbitrary existing pattern, validate general later
hierarchy-piece sources, or establish bounded hierarchy depth. Those remain
obligations for the original quasi-polynomial recognition goal. Ordinary
recognition is not yet changed.

\subsection{Source geometry and independent cut evidence}

The retained audit checks 1,755 finite cases: all 1,635 supplied pilot
surfaces, 48 mixed source cases, 36 analytic examples and 36 relabellings.
Every global genus inventory and cut-piece boundary count/total genus
agrees with Regina's actual three-dimensional cut. Forty-five cases contain
a piece with several boundary components. The literal cell audit checks
340,202 signed edge incidences, closed polygon cycles and connected vertex
links. The oracle uses \code{cutAlong} and each actual boundary component's
Euler characteristic\cite{regina-cut-along,regina-boundary-components}.

Ten large binary cases include parallel meridians, vertex links and odd/even
M\"obius multiples. A 16,385-bit multiplicity retains twelve polygon families
and 33 edge families in the one-tetrahedron meridian source, with 16 boundary
discovery cycles and 20 joint cycles. These numbers describe that supplied
family; the general bound remains polynomial in encoded source size.
All 110 saved source proofs independently replay with discovery, forward
transversal compilation, weighted producers and the cell producer disabled.
The source/evidence manifest authenticates 677 hashes. All ten focused
test methods pass, including mutations of endpoints, orientations, chamber
carriers, representative intervals and topology summaries, and shared caps.
The full maintained suite passes all 1,589 tests with Regina available and
no skips; its complete log is retained in \code{data/cut-boundary-tests.txt}.

The data are \code{data/cut-boundary-audit.json}. Reproduce the source audit
from \code{fast/} with:
\begin{quote}\small
\code{python -B -m normal\_orbit\_research.cut\_boundary --output FILE}
\end{quote}
The oracle is optional and requires the project's Regina extra; the native
constructor and independent checker use only the Python standard library.
The retained run uses the Regina Python distribution 7.4.1, whose engine
version string is 7.4, on Python 3.13.14.
